\subsection*{Discrete-Time Cross-Immunity Model}

The connectivity model assumes that all contacts lead to infection in
the same compartment. An alternative formulation applies to
cross-immunity between strains or pathogens, where exposure to a
related strain can confer immunity without causing infection by that
strain~\citep{gog2002}. We consider uni-directional cross-immunity:
exposure to strain $j$ confers immunity to strain $i$, but not vice
versa. The continuous-time equations are:
\begin{align*}
  \dot{S} &= -\beta(t) S \circ (I_d + M) I \\
  \dot{I} &= \beta(t) S\circ I - \nu I\\
  \dot{R} &= \beta(t) S \circ MI + \nu I,
\end{align*}
where the unidirectional cross-immunization matrix is
\begin{equation*}
M =
\begin{bmatrix}
  0 & \theta \\
  0 & 0
\end{bmatrix}
\end{equation*}
%% Anallogously, the force of infection in the discrete connectivity
%% model, we define the force of infection as 
%% \begin{equation*}
%% \lambda(t) = \beta(t) I(t)
%% \end{equation*}
The expected incidence is:
\begin{equation}\label{eq:cross_incidence}
\mu(t) = S(t) \left[1 - \exp(-\beta I) \right]
\end{equation}
The rest of the discrete-time dynamics arise naturally when we
carefully balance the terms to have a constant population:
\begin{align}\begin{split}\label{eq:discrete_cross}
S(t+1) &= S(t) \exp(-\beta (I_d + M)I) \\
I(t+1) &= I(t) \exp(-\gamma) + \mu(t) \\
R(t+1) &= R(t) + I(t) [1 - \exp(-\gamma)] + S(t)\exp(-\beta I(t))[1-\exp(-\beta MI(t))]
\end{split}\end{align}

\subsection*{Sensitivity equations}
The incidence sensitivity is:
\begin{align}\begin{split}\label{eq:sens_mu_cross}
    \frac{\partial \mu}{\partial \phi} %&= \frac{\partial}{\partial \phi}S[1-\exp(-\beta I) \\
    %
    %
    %
    &= \frac{\partial S}{\partial \phi} \left[1 - \exp(-\beta I)\right] +
    S \exp(-\beta I) \frac{\partial \beta I}{\partial \phi} \\
    %
    &=\frac{\partial S}{\partial \phi} \left[1 - \exp(-\beta I)\right] +
    \beta S \exp(-\beta I) \frac{\partial I}{\partial \phi}
\end{split}\end{align}
%% where
%% \begin{equation*}
%% \frac{\partial \lambda}{\partial \phi} = \beta
%% \frac{\partial I}{\partial \phi}
%% \end{equation*}
%
%
%% Note that $\theta$ does not appear in $\lambda_i$ or $\mu_i$
%% directly.  The parameter $\theta$ affects future incidence through
%% the depletion of susceptibles.
The state sensitivities propagate via:
\begin{align}\begin{split}\label{eq:sens_cross}
\frac{\partial S(t+1)}{\partial \phi} &=
\frac{\partial S}{\partial \phi} \exp(-\beta(I_d +M) I) - \beta S \exp(-\beta(I_d + M)I) \circ \left(
\frac{\partial I}{\partial \phi} + M\frac{\partial I}{\partial \phi} + \frac{\partial M}{\partial \phi}I \right) \\
\frac{\partial I(t+1)}{\partial \phi} &= \exp(-\gamma)
\frac{\partial I}{\partial \phi} + \frac{\partial \mu}{\partial \phi}
\end{split}\end{align}
%% nd
%% \begin{equation*}
%% \frac{\partial MI}{\partial \phi} = \beta \left(
%% \theta \frac{\partial I_j}{\partial \phi} + \delta_{\phi,\theta} I_j \right)
%% \end{equation*}
%% with $\delta_{\phi,\theta} = 1$ if $\phi = \theta$ and 0 otherwise.
%% The CRLB derivation proceeds as before, with
%% $\mathbf{J} = \mathbf{G}^T \mathbf{W} \mathbf{G}$.
